% Options for packages loaded elsewhere
\PassOptionsToPackage{unicode}{hyperref}
\PassOptionsToPackage{hyphens}{url}
\documentclass[
]{article}
\usepackage{xcolor}
\usepackage[margin=1in]{geometry}
\usepackage{amsmath,amssymb}
\setcounter{secnumdepth}{5}
\usepackage{iftex}
\ifPDFTeX
  \usepackage[T1]{fontenc}
  \usepackage[utf8]{inputenc}
  \usepackage{textcomp} % provide euro and other symbols
\else % if luatex or xetex
  \usepackage{unicode-math} % this also loads fontspec
  \defaultfontfeatures{Scale=MatchLowercase}
  \defaultfontfeatures[\rmfamily]{Ligatures=TeX,Scale=1}
\fi
\usepackage{lmodern}
\ifPDFTeX\else
  % xetex/luatex font selection
\fi
% Use upquote if available, for straight quotes in verbatim environments
\IfFileExists{upquote.sty}{\usepackage{upquote}}{}
\IfFileExists{microtype.sty}{% use microtype if available
  \usepackage[]{microtype}
  \UseMicrotypeSet[protrusion]{basicmath} % disable protrusion for tt fonts
}{}
\makeatletter
\@ifundefined{KOMAClassName}{% if non-KOMA class
  \IfFileExists{parskip.sty}{%
    \usepackage{parskip}
  }{% else
    \setlength{\parindent}{0pt}
    \setlength{\parskip}{6pt plus 2pt minus 1pt}}
}{% if KOMA class
  \KOMAoptions{parskip=half}}
\makeatother
\usepackage{color}
\usepackage{fancyvrb}
\newcommand{\VerbBar}{|}
\newcommand{\VERB}{\Verb[commandchars=\\\{\}]}
\DefineVerbatimEnvironment{Highlighting}{Verbatim}{commandchars=\\\{\}}
% Add ',fontsize=\small' for more characters per line
\usepackage{framed}
\definecolor{shadecolor}{RGB}{248,248,248}
\newenvironment{Shaded}{\begin{snugshade}}{\end{snugshade}}
\newcommand{\AlertTok}[1]{\textcolor[rgb]{0.94,0.16,0.16}{#1}}
\newcommand{\AnnotationTok}[1]{\textcolor[rgb]{0.56,0.35,0.01}{\textbf{\textit{#1}}}}
\newcommand{\AttributeTok}[1]{\textcolor[rgb]{0.13,0.29,0.53}{#1}}
\newcommand{\BaseNTok}[1]{\textcolor[rgb]{0.00,0.00,0.81}{#1}}
\newcommand{\BuiltInTok}[1]{#1}
\newcommand{\CharTok}[1]{\textcolor[rgb]{0.31,0.60,0.02}{#1}}
\newcommand{\CommentTok}[1]{\textcolor[rgb]{0.56,0.35,0.01}{\textit{#1}}}
\newcommand{\CommentVarTok}[1]{\textcolor[rgb]{0.56,0.35,0.01}{\textbf{\textit{#1}}}}
\newcommand{\ConstantTok}[1]{\textcolor[rgb]{0.56,0.35,0.01}{#1}}
\newcommand{\ControlFlowTok}[1]{\textcolor[rgb]{0.13,0.29,0.53}{\textbf{#1}}}
\newcommand{\DataTypeTok}[1]{\textcolor[rgb]{0.13,0.29,0.53}{#1}}
\newcommand{\DecValTok}[1]{\textcolor[rgb]{0.00,0.00,0.81}{#1}}
\newcommand{\DocumentationTok}[1]{\textcolor[rgb]{0.56,0.35,0.01}{\textbf{\textit{#1}}}}
\newcommand{\ErrorTok}[1]{\textcolor[rgb]{0.64,0.00,0.00}{\textbf{#1}}}
\newcommand{\ExtensionTok}[1]{#1}
\newcommand{\FloatTok}[1]{\textcolor[rgb]{0.00,0.00,0.81}{#1}}
\newcommand{\FunctionTok}[1]{\textcolor[rgb]{0.13,0.29,0.53}{\textbf{#1}}}
\newcommand{\ImportTok}[1]{#1}
\newcommand{\InformationTok}[1]{\textcolor[rgb]{0.56,0.35,0.01}{\textbf{\textit{#1}}}}
\newcommand{\KeywordTok}[1]{\textcolor[rgb]{0.13,0.29,0.53}{\textbf{#1}}}
\newcommand{\NormalTok}[1]{#1}
\newcommand{\OperatorTok}[1]{\textcolor[rgb]{0.81,0.36,0.00}{\textbf{#1}}}
\newcommand{\OtherTok}[1]{\textcolor[rgb]{0.56,0.35,0.01}{#1}}
\newcommand{\PreprocessorTok}[1]{\textcolor[rgb]{0.56,0.35,0.01}{\textit{#1}}}
\newcommand{\RegionMarkerTok}[1]{#1}
\newcommand{\SpecialCharTok}[1]{\textcolor[rgb]{0.81,0.36,0.00}{\textbf{#1}}}
\newcommand{\SpecialStringTok}[1]{\textcolor[rgb]{0.31,0.60,0.02}{#1}}
\newcommand{\StringTok}[1]{\textcolor[rgb]{0.31,0.60,0.02}{#1}}
\newcommand{\VariableTok}[1]{\textcolor[rgb]{0.00,0.00,0.00}{#1}}
\newcommand{\VerbatimStringTok}[1]{\textcolor[rgb]{0.31,0.60,0.02}{#1}}
\newcommand{\WarningTok}[1]{\textcolor[rgb]{0.56,0.35,0.01}{\textbf{\textit{#1}}}}
\usepackage{longtable,booktabs,array}
\usepackage{calc} % for calculating minipage widths
% Correct order of tables after \paragraph or \subparagraph
\usepackage{etoolbox}
\makeatletter
\patchcmd\longtable{\par}{\if@noskipsec\mbox{}\fi\par}{}{}
\makeatother
% Allow footnotes in longtable head/foot
\IfFileExists{footnotehyper.sty}{\usepackage{footnotehyper}}{\usepackage{footnote}}
\makesavenoteenv{longtable}
\usepackage{graphicx}
\makeatletter
\newsavebox\pandoc@box
\newcommand*\pandocbounded[1]{% scales image to fit in text height/width
  \sbox\pandoc@box{#1}%
  \Gscale@div\@tempa{\textheight}{\dimexpr\ht\pandoc@box+\dp\pandoc@box\relax}%
  \Gscale@div\@tempb{\linewidth}{\wd\pandoc@box}%
  \ifdim\@tempb\p@<\@tempa\p@\let\@tempa\@tempb\fi% select the smaller of both
  \ifdim\@tempa\p@<\p@\scalebox{\@tempa}{\usebox\pandoc@box}%
  \else\usebox{\pandoc@box}%
  \fi%
}
% Set default figure placement to htbp
\def\fps@figure{htbp}
\makeatother
% definitions for citeproc citations
\NewDocumentCommand\citeproctext{}{}
\NewDocumentCommand\citeproc{mm}{%
  \begingroup\def\citeproctext{#2}\cite{#1}\endgroup}
\makeatletter
 % allow citations to break across lines
 \let\@cite@ofmt\@firstofone
 % avoid brackets around text for \cite:
 \def\@biblabel#1{}
 \def\@cite#1#2{{#1\if@tempswa , #2\fi}}
\makeatother
\newlength{\cslhangindent}
\setlength{\cslhangindent}{1.5em}
\newlength{\csllabelwidth}
\setlength{\csllabelwidth}{3em}
\newenvironment{CSLReferences}[2] % #1 hanging-indent, #2 entry-spacing
 {\begin{list}{}{%
  \setlength{\itemindent}{0pt}
  \setlength{\leftmargin}{0pt}
  \setlength{\parsep}{0pt}
  % turn on hanging indent if param 1 is 1
  \ifodd #1
   \setlength{\leftmargin}{\cslhangindent}
   \setlength{\itemindent}{-1\cslhangindent}
  \fi
  % set entry spacing
  \setlength{\itemsep}{#2\baselineskip}}}
 {\end{list}}
\usepackage{calc}
\newcommand{\CSLBlock}[1]{\hfill\break\parbox[t]{\linewidth}{\strut\ignorespaces#1\strut}}
\newcommand{\CSLLeftMargin}[1]{\parbox[t]{\csllabelwidth}{\strut#1\strut}}
\newcommand{\CSLRightInline}[1]{\parbox[t]{\linewidth - \csllabelwidth}{\strut#1\strut}}
\newcommand{\CSLIndent}[1]{\hspace{\cslhangindent}#1}
\setlength{\emergencystretch}{3em} % prevent overfull lines
\providecommand{\tightlist}{%
  \setlength{\itemsep}{0pt}\setlength{\parskip}{0pt}}
\usepackage{float}
\usepackage{booktabs}
\usepackage{longtable}
\usepackage{array}
\usepackage{multirow}
\usepackage{wrapfig}
\usepackage{colortbl}
\usepackage{pdflscape}
\usepackage{tabu}
\usepackage{threeparttable}
\usepackage{threeparttablex}
\usepackage[normalem]{ulem}
\usepackage{makecell}
\usepackage{xcolor}
\usepackage{bookmark}
\IfFileExists{xurl.sty}{\usepackage{xurl}}{} % add URL line breaks if available
\urlstyle{same}
\hypersetup{
  pdftitle={Pràctica 3 : Cas-control per avaluar la relació entre exposició a fàrmacs antidepressius i fractura de maluc (Disseny aparellat).},
  pdfauthor={Iker Fernandez (1704341); Enola Garcia (1708253); Didac Campuzano (1703766); Marta Baurier (1668618); Jana Labrador (1704595)},
  hidelinks,
  pdfcreator={LaTeX via pandoc}}

\title{Pràctica 3 : Cas-control per avaluar la relació entre exposició a
fàrmacs antidepressius i fractura de maluc (Disseny aparellat).}
\author{Iker Fernandez (1704341) \and Enola Garcia (1708253) \and Didac
Campuzano (1703766) \and Marta Baurier (1668618) \and Jana Labrador
(1704595)}
\date{2026-04-22}

\begin{document}
\maketitle

L'objectiu d'aquest estudi és avaluar l'associació entre l'exposició a
fàrmacs antidepressius i el risc de fractura de maluc. En aquesta
anàlisi, s'han descartat els controls hospitalaris previs per utilitzar
nous controls procedents d'atenció primària, aparellats en una relació
1:4 per edat, sexe, centre i data de consulta. Això permet treballar amb
una població de referència més coherent, millorant la validesa externa
dels resultats

Creuem les dades de totes les bases de dades, el codi el podem veure a
l'annex\footnote{Annex \ref{anx:carrega}}

La Taula \ref{tab:tab1} mostra els estadístics descriptius de les
característiques de la població segons el grup d'estudi (controls i
casos). Per a les variables categòriques es presenten les freqüències
absolutes i els percentatges, mentre que per a les variables
quantitatives es mostren les mitjanes i les desviacions estàndard.

\begin{longtable}[t]{llccc}
\caption{\label{tab:unnamed-chunk-1}\label{tab:tab1}Característiques de la població}\\
\toprule
 & Control & Cas & P\_valor & \\
\midrule
n & 524 & 131 &  & \\
Sexe = M (\%) & 108 (20.6) & 27 (20.6) & 1.000 & \\
Edat (\%) &  &  & 1.000 & \\
\hspace{1em}51-65 & 28 ( 5.3) & 7 ( 5.3) &  & \\
\hspace{1em}66-81 & 172 (32.8) & 43 (32.8) &  & \\
\addlinespace
\hspace{1em}82-95 & 324 (61.8) & 81 (61.8) &  & \\
PES (mean (SD)) & 67.53 (13.04) & 64.84 (11.68) & 0.173 & \\
TALLA (mean (SD)) & 155.47 (8.02) & 152.36 (7.83) & 0.036 & \\
IMC (mean (SD)) & 27.76 (4.83) & 27.55 (4.90) & 0.809 & \\
Smoker (\%) &  &  & 0.001 & \\
\addlinespace
No Fumador & 502 (95.8) & 117 (89.3) &  & \\
\hspace{1em}Fumador Actiu & 16 ( 3.1) & 6 ( 4.6) &  & \\
\hspace{1em}Ex-Fumador & 6 ( 1.1) & 8 ( 6.1) &  & \\
\hspace{1em}ALCOHOL (\%) &  &  & 0.058 & \\
No consumeix & 465 (88.7) & 123 (93.9) &  & \\
\addlinespace
\hspace{1em}Poc consum & 56 (10.7) & 6 ( 4.6) &  & \\
\hspace{1em}Consum moderat & 3 ( 0.6) & 2 ( 1.5) &  & \\
\hspace{1em}Artritis\_reumatoide = 1 (\%) & 5 ( 1.0) & 1 ( 0.8) & 1.000 & \\
Fractura = 1 (\%) & 2 ( 0.4) & 55 (42.0) & <0.001 & \\
NFractura (\%) &  &  & <0.001 & \\
\addlinespace
0 & 522 (99.6) & 76 (58.0) &  & \\
\hspace{1em}1 & 2 ( 0.4) & 48 (36.6) &  & \\
\hspace{1em}>=2 & 0 ( 0.0) & 7 ( 5.3) &  & \\
\hspace{1em}Diabetis = 1 (\%) & 118 (22.5) & 45 (34.4) & 0.007 & \\
tipus\_diabetis (\%) &  &  & 0.007 & \\
\addlinespace
No diabetis & 406 (77.5) & 86 (65.6) &  & \\
\hspace{1em}Tipus 1 & 4 ( 0.8) & 0 ( 0.0) &  & \\
\hspace{1em}Tipus 2 & 114 (21.8) & 45 (34.4) &  & \\
\hspace{1em}Osteporosi = 1 (\%) & 13 ( 2.5) & 17 (13.0) & <0.001 & \\
Densitometries = 1 (\%) & 0 ( 0.0) & 8 ( 6.1) & <0.001 & \\
\addlinespace
neoplasia = 1 (\%) & 61 (11.6) & 12 ( 9.2) & 0.514 & \\
HiperTiroidisme = 1 (\%) & 3 ( 0.6) & 4 ( 3.1) & 0.046 & \\
Malnutricio = 1 (\%) & 1 ( 0.2) & 3 ( 2.3) & 0.033 & \\
Malabsorcio = 1 (\%) & 1 ( 0.2) & 0 ( 0.0) & 1.000 & \\
Malaltia\_Hep\_Cro = 1 (\%) & 7 ( 1.3) & 5 ( 3.8) & 0.126 & \\
\addlinespace
OthRiskFract = 1 (\%) & 75 (14.3) & 22 (16.8) & 0.564 & \\
CountActualGE4 = 1 (\%) & 222 (42.4) & 99 (75.6) & <0.001 & \\
\bottomrule
\end{longtable}

Cal destacar que diverses variables presenten una proporció elevada de
valors perduts (NA). Tot i que aquestes variables s'han inclòs en
l'anàlisi descriptiva, els seus resultats s'han d'interpretar amb
cautela, ja que aquest alt contingut de dades mancants pot limitar la
seva fiabilitat i representativitat. Per aquest motiu, en la
interpretació global no es dona el mateix pes a totes les variables.

En general, la majoria de variables presenten distribucions similars
entre els grups de casos i controls, sense diferències estadísticament
significatives. Això s'observa, per exemple, en variables com el sexe,
l'edat, el MatchCC, l'IMC o la presència de neoplàsia, on els valors són
molt semblants entre ambdós grups.

Tanmateix, s'observen diferències estadísticament significatives en
algunes variables com el servei d'alta, la talla, el consum de tabac
(Smoker), la presència de fractura i el nombre de fractures, la diabetis
i el tipus de diabetis, l'osteoporosi, la realització de densitometries,
l'hipertiroïdisme, la malnutrició i la variable CountActualGE4.

Pel que fa als mètodes estadístics utilitzats, per a les variables
categòriques s'han aplicat proves de comparació de proporcions,
principalment el test de chi-quadrat o el test exacte de Fisher quan les
freqüències eren baixes. Per a les variables quantitatives s'ha
utilitzat el test no paramètric de Wilcoxon (Mann-Whitney), adequat en
situacions on no es pot assumir normalitat en la distribució de les
dades.

Un element destacable és que variables directament relacionades amb el
resultat, com la fractura i el nombre de fractures, presenten
diferències molt marcades entre casos i controls, fet esperable atès que
formen part de la definició dels grups. També s'observa una major
prevalença d'algunes condicions clíniques com la diabetis o
l'osteoporosi en els casos, la qual cosa podria suggerir una possible
associació amb el risc de fractura.

No obstant això, aquesta anàlisi és purament descriptiva i no té en
compte possibles factors de confusió. A més, la presència de dades
mancants en algunes variables pot introduir incertesa en les
estimacions. Per tant, el fet que una variable presenti una major
proporció en un grup no implica necessàriament una relació causal. Per
aquest motiu, en les anàlisis posteriors es durà a terme una regressió
logística, tant en forma crua com ajustada per covariables, per tal
d'avaluar aquestes associacions de manera més robusta.

\clearpage

\section{Analitzar la exposició als fàrmacs i els riscs associats de
fractura}\label{analitzar-la-exposiciuxf3-als-fuxe0rmacs-i-els-riscs-associats-de-fractura}

\subsection{Anàlisi descriptiva d'exposició en els grups de cas i
control}\label{anuxe0lisi-descriptiva-dexposiciuxf3-en-els-grups-de-cas-i-control}

La taula \ref{tab:descrip} mostra la distribució dels diferents fàrmacs
i grups terapèutics entre casos i controls, permetent comparar el
percentatge d'exposició en ambdós grups.

\begin{longtable}[]{@{}llccc@{}}
\caption{\label{tab:descrip}Comparació de Medicaments i
teràpies}\tabularnewline
\toprule\noalign{}
& Control & Cas & P\_valor & \\
\midrule\noalign{}
\endfirsthead
\toprule\noalign{}
& Control & Cas & P\_valor & \\
\midrule\noalign{}
\endhead
\bottomrule\noalign{}
\endlastfoot
n & 524 & 131 & & \\
CortInh = 1 (\%) & 25 ( 4.8) & 15 (11.5) & 0.008 & \\
CortSist = 1 (\%) & 6 ( 1.1) & 5 ( 3.8) & 0.080 & \\
CortSistIniBf3m = 1 (\%) & 6 ( 1.1) & 3 ( 2.3) & 0.557 & \\
CortSistExpLt3m = 1 (\%) & 6 ( 1.1) & 3 ( 2.3) & 0.557 & \\
CortInhIniBf3m = 1 (\%) & 20 ( 3.8) & 10 ( 7.6) & 0.102 & \\
CortInhExpLt3m = 1 (\%) & 20 ( 3.8) & 10 ( 7.6) & 0.102 & \\
ADO\_no\_glitazona = 1 (\%) & 47 ( 9.0) & 26 (19.8) & 0.001 & \\
Bisf = 1 (\%) & 15 ( 2.9) & 17 (13.0) & \textless0.001 & \\
ADepreISRS = 1 (\%) & 40 ( 7.6) & 36 (27.5) & \textless0.001 & \\
ADepreNoISRS = 1 (\%) & 48 ( 9.2) & 30 (22.9) & \textless0.001 & \\
insulina = 1 (\%) & 20 ( 3.8) & 10 ( 7.6) & 0.102 & \\
H\_SN = 1 (\%) & 30 ( 5.7) & 15 (11.5) & 0.034 & \\
N\_SN = 1 (\%) & 223 (42.6) & 99 (75.6) & \textless0.001 & \\
H01\_SN = 1 (\%) & 1 ( 0.2) & 0 ( 0.0) & 1.000 & \\
H02\_SN = 1 (\%) & 7 ( 1.3) & 5 ( 3.8) & 0.126 & \\
N06\_SN = 1 (\%) & 83 (15.8) & 55 (42.0) & \textless0.001 & \\
N06A\_SN = 1 (\%) & 66 (12.6) & 50 (38.2) & \textless0.001 & \\
N06AA\_SN = 1 (\%) & 9 ( 1.7) & 4 ( 3.1) & 0.528 & \\
N06AB\_SN = 1 (\%) & 41 ( 7.8) & 36 (27.5) & \textless0.001 & \\
N06AX\_SN = 1 (\%) & 21 ( 4.0) & 15 (11.5) & 0.002 & \\
R03BA\_SN = 1 (\%) & 7 ( 1.3) & 2 ( 1.5) & 1.000 & \\
\end{longtable}

Els resultats de l'anàlisi descriptiva mostren diferències
significatives entre casos i controls en diversos grups farmacològics.
En particular, destaca una major prevalença d'ús d'antidepressius en el
grup de casos, tant per als ISRS com per als no ISRS. Els ISRS presenten
una prevalença del 27,5\% en casos enfront del 7,6\% en controls, i els
antidepressius no ISRS del 22,9\% davant del 9,2\%, amb p-valors
significatius en tots dos casos. Aquest mateix patró també s'observa en
diverses categories del sistema nerviós, com N06, N06A i N06AB.

També s'observen diferències significatives en altres grups terapèutics,
com els bifosfonats o els antidiabètics orals, fet que suggereix una
major càrrega terapèutica en el grup de casos.

En aquest sentit, (\citeproc{ref-coupland2011antidepressant}{Coupland et
al. 2011}) assenyalen que l'associació entre antidepressius i
esdeveniments adversos pot estar parcialment explicada per les
característiques dels pacients tractats. Aquest fet encaixa amb els
nostres resultats, perquè a la taula veiem que els casos prenen més
antidepressius que els controls, sobretot ISRS i antidepressius no ISRS,
tal com hem esmentat.

Per altra banda, veiem que els casos prenen més fàrmacs del sistema
nerviós, sobretot antidepressius. Segons la
(\citeproc{ref-brotomodulo}{Broto, Brumós, and Montañá, n.d.}), aquest
tipus de pacients sovint presenten polimedicació, i això augmenta el
risc d'efectes adversos i d'interaccions entre medicaments; és per això
que cal interpretar-ho amb prudència.

Per tant, serà necessari realitzar una anàlisi multivariant per
determinar si aquestes associacions es mantenen després d'ajustar per
possibles factors confusors.

\clearpage

\subsection{Anàlisi inferencial cru i
ajustat}\label{anuxe0lisi-inferencial-cru-i-ajustat}

A continuació, utilitzarem un model de Regressió Logística Condicional,
fent servir la funció \texttt{clogit} del paquet \texttt{survival}.
S'utilitza aquest tipus de model degut a que parlem d'un estudi
Cas-Control, on les dades son pariades, en aquest cas, 1:4, és a dir,
cada cas té 4 controls, aparellats en edat i sexe.

I per tant, les observacions dins de cada conjunt de 4 són dependents,
això fa que la regressió logística convencional no sigui adequada en
aquest cas.

\subsubsection{Anàlisi logística
``crua''}\label{anuxe0lisi-loguxedstica-crua}

La taula \ref{tab:logcru} presenta els resultats de la regressió
logística crua, que permet veure l'associació directa entre l'exposició
als diferents fàrmacs i el risc de fractura de maluc, sense ajustar per
possibles factors de confusió.

\begin{longtable}[]{@{}lcccl@{}}
\caption{\label{tab:logcru}Anàlisi per regressió logística crua de les
variables}\tabularnewline
\toprule\noalign{}
Variable & OR\_Cru & IC\_Inferior & IC\_Superior & P\_valor \\
\midrule\noalign{}
\endfirsthead
\toprule\noalign{}
Variable & OR\_Cru & IC\_Inferior & IC\_Superior & P\_valor \\
\midrule\noalign{}
\endhead
\bottomrule\noalign{}
\endlastfoot
CortInh & 2.40 & 1.27 & 4.55 & 0.007 \\
CortSist & 3.33 & 1.02 & 10.92 & 0.047 \\
CortSistIniBf3m & 2.00 & 0.50 & 8.00 & 0.327 \\
CortSistExpLt3m & 2.00 & 0.50 & 8.00 & 0.327 \\
CortInhIniBf3m & 2.00 & 0.94 & 4.27 & 0.074 \\
CortInhExpLt3m & 2.00 & 0.94 & 4.27 & 0.074 \\
ADO\_no\_glitazona & 2.56 & 1.50 & 4.37 & 0.001 \\
Bisf & 4.94 & 2.39 & 10.21 & 0.000 \\
ADepreISRS & 5.15 & 2.97 & 8.93 & 0.000 \\
ADepreNoISRS & 3.13 & 1.85 & 5.30 & 0.000 \\
insulina & 2.08 & 0.95 & 4.56 & 0.066 \\
H\_SN & 2.11 & 1.10 & 4.03 & 0.024 \\
N\_SN & 4.32 & 2.77 & 6.74 & 0.000 \\
H01\_SN & 0.00 & 0.00 & Inf & 0.997 \\
H02\_SN & 2.86 & 0.91 & 9.00 & 0.073 \\
N06\_SN & 4.16 & 2.65 & 6.52 & 0.000 \\
N06A\_SN & 4.58 & 2.87 & 7.31 & 0.000 \\
N06AA\_SN & 1.83 & 0.54 & 6.14 & 0.329 \\
N06AB\_SN & 4.93 & 2.87 & 8.49 & 0.000 \\
N06AX\_SN & 2.98 & 1.51 & 5.90 & 0.002 \\
R03BA\_SN & 1.14 & 0.24 & 5.50 & 0.868 \\
\end{longtable}

Els resultats mostren una associació positiva i estadísticament
significativa entre diversos fàrmacs i el risc de fractura de maluc. A
diferència d'altres dissenys amb controls hospitalaris, l'ús de controls
de població general podria contribuir a una major diferència en el
perfil d'exposició entre casos i controls.

Els resultats més destacats es troben en els \textbf{antidepressius}
(\texttt{N06\_SN}), amb una OR crua de 4.16 (IC: 2.65, 6.52). Dins
d'aquest grup, els ISRS (\texttt{N06AB\_SN}) presenten el risc més
eleva, amb una OR de 4.93, en línia amb el que descriu l'article de
referència (\citeproc{ref-gorgas2021}{Gorgas et al. 2021}), que
siggereix tant un augmet del risc de caigudes com possibles efectes
sobre el metabolisme ossi.

També s'observen riscos molt elevats en els \textbf{bisfosfonats}
(OR=4.94). Aquestes associacions s'han d'interpretar amb precaució, ja
que poden reflectir confusió per indicació: els pacients exposats són
probablement més fràgils o amb més patologies de base (com osteoporosi o
diabetis) que els controls. Aquesta gran diferència basal explica per
què les OR crues són tan elevades en comparació amb estudis on els
controls també estan malalts.

D'altra banda, algunes variables no mostren associacions estadísticament
significatives, com la insulina (p=0.066) o determinats subgrups
terapèutics. Finalment, la presència d'intervals de confiança molt
amplis o valors extrems, com en H01\_SN, indica possibles problemes de
baixa freqüència o separació quasi perfecta, fet que limita la
fiabilitat d'aquestes estimacions.

\subsubsection{Anàlisi logística
ajustada}\label{anuxe0lisi-loguxedstica-ajustada}

L'anàlisi crua detecta un increment de risc generalitzat, però la
magnitud d'aquestes OR suggereix la presència d'un fort biaix de
confusió per indicació. Per tant, l'anàlisi ajustada que es veu a
continuació, és imprescindible per intentar aïllar l'efecte real dels
antidepressius de la fragilitat general del pacient.

Per tal de fer el millor ajust possible, observem la taula
\ref{tab:tab1} per observar possibles variables confusores, i tenim que
aquestes són:

\begin{itemize}
\item
  \textbf{Variables confusores retingudes}
\item
  \texttt{Smoker} i \texttt{Alcohol}: S'ha decidit incloure aquestes
  variables en l'anàlisi ajustat ja que són factors d'estil de vida que
  afecten directament la salut òssia i poden actuar com a variables de
  confusió en l'associació amb els medicaments
  (\citeproc{ref-leal2022description}{\textbf{leal2022description?}})
  (\citeproc{ref-botaya2018osteoporosis}{\textbf{botaya2018osteoporosis?}}).
  A la descripció basal (\ref{tab:tab1}), s'observen diferències
  estadísticament significatives en el consum de tabac (p=0.001), amb
  una major proporció de fumadors actius i ex-fumadorsa l grup dels
  casos. Pel que fa a l'alcohol (p=0,058), tot i trobar-se en el límit
  de la significació, s'ha optat per mantenir-la al model ajustat per la
  seva rellevància clínica i bibliogràfica en el risc de fractures i
  caigudes, a més, aquesta diferéncia es troba als controls en l'àmbit
  de poc consum.
\item
  \texttt{Diabetis}, \texttt{CountActualGE4}: De manera similar, es va
  avaluar l'impacte de prendre 4 o més fàrmacs de risc (CountActualGE4),
  ja que pot actuar com a indicador de polimedicació i de major càrrega
  de comorbiditat
  (\citeproc{ref-bonaga2016polifarmacia}{\textbf{bonaga2016polifarmacia?}}).
  En la descripció basal (\ref{tab:tab1}), aquesta variable mostra
  diferències estadísticament significatives entre casos i controls(
  p\textless0.001), amb una major proporció de pacients polimedicats en
  el grup de casos (75.6\% vs 42.4\%).
\end{itemize}

Pel que fa a la diabetis, també s'observen diferències significatives
(p=0.007), amb una major prevalença en els casos. Donada la seva relació
amb complicacions metabòliques i possibles efectes sobre el risc de
fractura, així com l'ús de tractaments com la insulina i altres
antidiabètics
(\citeproc{ref-sedlinskydiabetes}{\textbf{sedlinskydiabetes?}}), es
considera una variable rellevant a incloure com a covariable en el model
ajustat.

\begin{itemize}
\item
  \texttt{Osteoporosi} i \texttt{Fractura}: Compleixen un doble criteri.
  D'una banda, presenten una forta evidència clínica com a factors de
  risc directes per a futures fractures. L'osteporosi és una malaltia
  que debilita els ossos i augmenta el risc de fractures i, Fractura
  indica si han hagut fractures prèvies
  (\citeproc{ref-campos2019osteoporosis}{\textbf{campos2019osteoporosis?}}).
  D'altra banda, l'anàlisi de comparació de grups confirma una
  diferència estadísticament significativa en la nostra mostra d'estudi.
\item
  \textbf{Variables avaluades i excloses}:
\item
  \texttt{Talla}, \texttt{HiperTiroidisme} i \texttt{Malnutricio}: La
  variable talla, ha estat exclosa degut a la presència de NA's a la
  base de dades, cosa que fa que els resultats s'esbiaixin. Per la part
  de l' HiperTiroidisme i la Malnutricio, tot i que hii ha diferencies
  significatives, s'ha optat per no incloure-les ja que aquestes tenen
  una baixa freqüència d'aquestes als casos i controls, la introducció
  d'aquestes variables podria generar problemes d'inestabilitat
  numèrica, a més, s'ha considerat que el seu potencial com a confusor,
  pot estar bastant limitat.
\end{itemize}

Per tant, un cop s'han triat els confusors, fem el model:

\begin{longtable}[]{@{}lcccc@{}}
\caption{\label{tab:logajust}Anàlisi per regressió logística ajustada
per covariables}\tabularnewline
\toprule\noalign{}
Variable & OR\_Ajustat & IC\_Inferior & IC\_Superior & P\_valor \\
\midrule\noalign{}
\endfirsthead
\toprule\noalign{}
Variable & OR\_Ajustat & IC\_Inferior & IC\_Superior & P\_valor \\
\midrule\noalign{}
\endhead
\bottomrule\noalign{}
\endlastfoot
CortInh & 1.78 & 0.69 & 4.61 & 0.233 \\
CortSist & 3.94 & 0.83 & 18.74 & 0.085 \\
CortSistIniBf3m & 3.09 & 0.58 & 16.50 & 0.188 \\
CortSistExpLt3m & 3.09 & 0.58 & 16.50 & 0.188 \\
CortInhIniBf3m & 1.60 & 0.55 & 4.66 & 0.393 \\
CortInhExpLt3m & 1.60 & 0.55 & 4.66 & 0.393 \\
ADO\_no\_glitazona & 1.40 & 0.53 & 3.72 & 0.495 \\
Bisf & 1.64 & 0.42 & 6.41 & 0.477 \\
ADepreISRS & 3.72 & 1.62 & 8.54 & 0.002 \\
ADepreNoISRS & 1.62 & 0.72 & 3.67 & 0.245 \\
insulina & 1.29 & 0.40 & 4.11 & 0.668 \\
H\_SN & 1.21 & 0.48 & 3.07 & 0.681 \\
N\_SN & 2.86 & 1.31 & 6.25 & 0.008 \\
H01\_SN & 0.00 & 0.00 & Inf & 0.997 \\
H02\_SN & 3.00 & 0.70 & 12.81 & 0.138 \\
N06\_SN & 2.34 & 1.17 & 4.70 & 0.017 \\
N06A\_SN & 2.36 & 1.16 & 4.81 & 0.018 \\
N06AA\_SN & 0.92 & 0.12 & 6.99 & 0.935 \\
N06AB\_SN & 3.47 & 1.53 & 7.84 & 0.003 \\
N06AX\_SN & 0.79 & 0.24 & 2.59 & 0.696 \\
R03BA\_SN & 1.69 & 0.25 & 11.32 & 0.590 \\
\end{longtable}

Els resultats de la regressió logística ajustada mostren una reducció
general de les odds ratio respecte a l'anàlisi crua, fet que suggereix
la presència de confusió per indicació en molts dels fàrmacs analitzats.
Després d'ajustar per covariables com diabetis, alcohol, tabaquisme i
osteoporosi, la majoria d'associacions perden significació estadística.

Els resultats més destacats es mantenen en els antidepressius,
especialment dins del grup N06, amb una OR ajustada de 2.34 (IC: 1.17,
4.70). Dins d'aquest grup, els ISRS (N06AB) continuen presentant una
associació significativa (OR=3.47; IC: 1.53, 7.84), fet que reforça la
hipòtesi d'un possible efecte propi d'aquests fàrmacs més enllà del
perfil basal dels pacients.

En canvi, altres fàrmacs que mostraven associacions elevades en
l'anàlisi crua, com els bisfosfonats o la insulina, deixen de ser
estadísticament significatius després de l'ajust. Això suggereix que les
associacions observades inicialment podrien estar influïdes per la
presència de factors de confusió.

\subsection{Comparació amb l'article}\label{comparaciuxf3-amb-larticle}

Els resultats obtinguts en aquest estudi són, en general, similars amb
els descrits a l'article de referència (\citeproc{ref-gorgas2021}{Gorgas
et al. 2021}), tant en l'anàlisi crua com en l'ajustada.

Pel que fa a l'anàlisi crua, les magnituds de les odds ratio són molt
similars entre ambdós estudis. En particular, els antidepressius,
especialment els ISRS (N06AB), mostren una OR elevada tant en el nostre
estudi (OR =4.93) com en l'article (OR=4.89), indicant una forta
associació entre exposició i risc de fractura abans d'ajustar per
factors de confusió.

En l'anàlisi ajustada, es manté aquesta coherència. Els ISRS continuen
mostrant una associació positiva i estadísticament significativa en
ambdós estudis, amb una OR de 3.47 en els nostres resultats i de 3.52 en
l'article.

Altres antidepressius no ISRS perden significació estadística després de
l'ajust en ambdós estudis, suggerint que l'associació observada en
l'anàlisi crua està influïda per factors de confusió.

Les petites diferències observades entre els resultats poden deure's a
variacions en les covariables incloses en el model, cosa que podria
explicar les discrepàncies observades entre ambdós estudis.

\clearpage

\section{Estadístics descriptius, mesures de risc,
inferencies\ldots{}}\label{estaduxedstics-descriptius-mesures-de-risc-inferencies}

\begin{longtable}[]{@{}
  >{\raggedright\arraybackslash}p{(\linewidth - 8\tabcolsep) * \real{0.1635}}
  >{\centering\arraybackslash}p{(\linewidth - 8\tabcolsep) * \real{0.1538}}
  >{\centering\arraybackslash}p{(\linewidth - 8\tabcolsep) * \real{0.1250}}
  >{\centering\arraybackslash}p{(\linewidth - 8\tabcolsep) * \real{0.2692}}
  >{\centering\arraybackslash}p{(\linewidth - 8\tabcolsep) * \real{0.2885}}@{}}
\caption{Exposició a fàrmacs i riscos associats}\tabularnewline
\toprule\noalign{}
\begin{minipage}[b]{\linewidth}\raggedright
Exposure
\end{minipage} & \begin{minipage}[b]{\linewidth}\centering
Controls n (\%)
\end{minipage} & \begin{minipage}[b]{\linewidth}\centering
Casos n (\%)
\end{minipage} & \begin{minipage}[b]{\linewidth}\centering
OR crua (IC95\%); p valor
\end{minipage} & \begin{minipage}[b]{\linewidth}\centering
OR ajustada (IC95\%); p valor
\end{minipage} \\
\midrule\noalign{}
\endfirsthead
\toprule\noalign{}
\begin{minipage}[b]{\linewidth}\raggedright
Exposure
\end{minipage} & \begin{minipage}[b]{\linewidth}\centering
Controls n (\%)
\end{minipage} & \begin{minipage}[b]{\linewidth}\centering
Casos n (\%)
\end{minipage} & \begin{minipage}[b]{\linewidth}\centering
OR crua (IC95\%); p valor
\end{minipage} & \begin{minipage}[b]{\linewidth}\centering
OR ajustada (IC95\%); p valor
\end{minipage} \\
\midrule\noalign{}
\endhead
\bottomrule\noalign{}
\endlastfoot
CortInh & 25 ( 4.8) & 15 (11.5) & 2.4 (1.27,4.55); p=0.007 & 1.78
(0.69,4.61); p=0.233 \\
CortSist & 6 ( 1.1) & 5 ( 3.8) & 3.33 (1.02,10.92); p=0.047 & 3.94
(0.83,18.74); p=0.085 \\
CortSistIniBf3m & 6 ( 1.1) & 3 ( 2.3) & 2 (0.5,8.00); p=0.327 & 3.09
(0.58,16.50); p=0.188 \\
CortSistExpLt3m & 6 ( 1.1) & 3 ( 2.3) & 2 (0.5,8.00); p=0.327 & 3.09
(0.58,16.50); p=0.188 \\
CortInhIniBf3m & 20 ( 3.8) & 10 ( 7.6) & 2 (0.94,4.27); p=0.074 & 1.6
(0.55,4.66); p=0.393 \\
CortInhExpLt3m & 20 ( 3.8) & 10 ( 7.6) & 2 (0.94,4.27); p=0.074 & 1.6
(0.55,4.66); p=0.393 \\
ADO\_no\_glitazona & 47 ( 9.0) & 26 (19.8) & 2.56 (1.5,4.37); p=0.001 &
1.4 (0.53,3.72); p=0.495 \\
Bisf & 15 ( 2.9) & 17 (13.0) & 4.94 (2.39,10.21); p=0 & 1.64
(0.42,6.41); p=0.477 \\
ADepreISRS & 40 ( 7.6) & 36 (27.5) & 5.15 (2.97,8.93); p=0 & 3.72
(1.62,8.54); p=0.002 \\
ADepreNoISRS & 48 ( 9.2) & 30 (22.9) & 3.13 (1.85,5.30); p=0 & 1.62
(0.72,3.67); p=0.245 \\
insulina & 20 ( 3.8) & 10 ( 7.6) & 2.08 (0.95,4.56); p=0.066 & 1.29
(0.4,4.11); p=0.668 \\
H\_SN & 30 ( 5.7) & 15 (11.5) & 2.11 (1.1,4.03); p=0.024 & 1.21
(0.48,3.07); p=0.681 \\
N\_SN & 223 (42.6) & 99 (75.6) & 4.32 (2.77,6.74); p=0 & 2.86
(1.31,6.25); p=0.008 \\
H01\_SN & 1 ( 0.2) & 0 ( 0.0) & 0 (0,Inf); p=0.997 & 0 (0,Inf);
p=0.997 \\
H02\_SN & 7 ( 1.3) & 5 ( 3.8) & 2.86 (0.91,9.00); p=0.073 & 3
(0.7,12.81); p=0.138 \\
N06\_SN & 83 (15.8) & 55 (42.0) & 4.16 (2.65,6.52); p=0 & 2.34
(1.17,4.70); p=0.017 \\
N06A\_SN & 66 (12.6) & 50 (38.2) & 4.58 (2.87,7.31); p=0 & 2.36
(1.16,4.81); p=0.018 \\
N06AA\_SN & 9 ( 1.7) & 4 ( 3.1) & 1.83 (0.54,6.14); p=0.329 & 0.92
(0.12,6.99); p=0.935 \\
N06AB\_SN & 41 ( 7.8) & 36 (27.5) & 4.93 (2.87,8.49); p=0 & 3.47
(1.53,7.84); p=0.003 \\
N06AX\_SN & 21 ( 4.0) & 15 (11.5) & 2.98 (1.51,5.90); p=0.002 & 0.79
(0.24,2.59); p=0.696 \\
R03BA\_SN & 7 ( 1.3) & 2 ( 1.5) & 1.14 (0.24,5.50); p=0.868 & 1.69
(0.25,11.32); p=0.59 \\
\end{longtable}

\section{Factors de confusió}\label{factors-de-confusiuxf3}

Els factors de confusió són variables que fan que una relació aparent
entre un fàrmac i el risc de fractura pugui estar influïda, en realitat,
per característiques dels pacients i no únicament per l'efecte del
tractament.

En aquest estudi, la Taula 1 i la Taula 2 ja mostren diferències entre
casos i controls en diverses variables clíniques i d'estil de vida.
També s'observa que diversos medicaments són més freqüents en el grup de
casos, especialment antidepressius, bisfosfonats i antidiabètics. Tot i
això, aquestes diferències descriptives no permeten concloure una
relació causal directa.

La comparació entre la Taula 3 (anàlisi crua) i la Taula 4 (anàlisi
ajustada) permet valorar millor la presència de confusió. Quan una odds
ratio disminueix de manera notable després de l'ajust o deixa de ser
significativa, això suggereix que part de l'associació inicial estava
influïda per factors de confusió.

En aquest cas, els bisfosfonats passen d'una OR crua elevada a una OR
ajustada molt menor i no significativa. Aquest patró és coherent amb
confusió per indicació, ja que els pacients que reben aquest tractament
presenten amb més freqüència osteoporosi o major risc ossi previ. Un
comportament similar s'observa en la insulina i altres tractaments
metabòlics, on la càrrega de comorbiditat podria explicar part de
l'excés de risc observat en l'anàlisi no ajustada.

Cal destacar també que la pèrdua de significació després de l'ajust no
vol dir necessàriament que no hi hagi efecte, sinó que l'efecte estimat
pot ser més petit, menys precís o ambdues coses alhora. Que els
intervals de confiança siguin més amplis després d'afegir covariables és
habitual en models amb diverses variables, especialment quan la mostra
no és molt gran.

Finalment, no es pot descartar la presència de confusió residual per
factors que no estan disponibles a la base de dades, com antecedents de
caigudes, estat funcional, gravetat clínica, dosi acumulada o adherència
al tractament. Per això, els resultats s'han d'interpretar com
associacions ajustades i no com una prova causal definitiva. Tot i així,
l'ús de regressió logística condicional en un disseny aparellat continua
sent una opció adequada per obtenir resultats més fiables en aquest
context.

\begin{longtable}[]{@{}
  >{\raggedright\arraybackslash}p{(\linewidth - 10\tabcolsep) * \real{0.1809}}
  >{\centering\arraybackslash}p{(\linewidth - 10\tabcolsep) * \real{0.0851}}
  >{\centering\arraybackslash}p{(\linewidth - 10\tabcolsep) * \real{0.1277}}
  >{\centering\arraybackslash}p{(\linewidth - 10\tabcolsep) * \real{0.2021}}
  >{\centering\arraybackslash}p{(\linewidth - 10\tabcolsep) * \real{0.1277}}
  >{\raggedright\arraybackslash}p{(\linewidth - 10\tabcolsep) * \real{0.2766}}@{}}
\caption{Avaluació de factors de confusió mitjançant el canvi relatiu de
l'OR}\tabularnewline
\toprule\noalign{}
\begin{minipage}[b]{\linewidth}\raggedright
Variable
\end{minipage} & \begin{minipage}[b]{\linewidth}\centering
OR Cru
\end{minipage} & \begin{minipage}[b]{\linewidth}\centering
OR Ajustat
\end{minipage} & \begin{minipage}[b]{\linewidth}\centering
Canvi relatiu (\%)
\end{minipage} & \begin{minipage}[b]{\linewidth}\centering
Confusora?
\end{minipage} & \begin{minipage}[b]{\linewidth}\raggedright
Direcció
\end{minipage} \\
\midrule\noalign{}
\endfirsthead
\toprule\noalign{}
\begin{minipage}[b]{\linewidth}\raggedright
Variable
\end{minipage} & \begin{minipage}[b]{\linewidth}\centering
OR Cru
\end{minipage} & \begin{minipage}[b]{\linewidth}\centering
OR Ajustat
\end{minipage} & \begin{minipage}[b]{\linewidth}\centering
Canvi relatiu (\%)
\end{minipage} & \begin{minipage}[b]{\linewidth}\centering
Confusora?
\end{minipage} & \begin{minipage}[b]{\linewidth}\raggedright
Direcció
\end{minipage} \\
\midrule\noalign{}
\endhead
\bottomrule\noalign{}
\endlastfoot
CortInh & 2.40 & 1.80 & 33.3 & Sí (+10\%) & Positiva (sobreestimació) \\
CortSist & 3.33 & 3.93 & -15.3 & Sí (+10\%) & Negativa (subestimació) \\
CortSistIniBf3m & 2.00 & 3.05 & -34.4 & Sí (+10\%) & Negativa
(subestimació) \\
CortSistExpLt3m & 2.00 & 3.05 & -34.4 & Sí (+10\%) & Negativa
(subestimació) \\
CortInhIniBf3m & 2.00 & 1.61 & 24.2 & Sí (+10\%) & Positiva
(sobreestimació) \\
CortInhExpLt3m & 2.00 & 1.61 & 24.2 & Sí (+10\%) & Positiva
(sobreestimació) \\
ADO\_no\_glitazona & 2.56 & 1.40 & 82.9 & Sí (+10\%) & Positiva
(sobreestimació) \\
Bisf & 4.94 & 1.13 & 337.2 & Sí (+10\%) & Positiva (sobreestimació) \\
ADepreISRS & 5.15 & 2.83 & 82.0 & Sí (+10\%) & Positiva
(sobreestimació) \\
ADepreNoISRS & 3.13 & 1.46 & 114.4 & Sí (+10\%) & Positiva
(sobreestimació) \\
insulina & 2.08 & 1.18 & 76.3 & Sí (+10\%) & Positiva
(sobreestimació) \\
H\_SN & 2.11 & 1.52 & 38.8 & Sí (+10\%) & Positiva (sobreestimació) \\
N\_SN & 4.32 & 2.34 & 84.6 & Sí (+10\%) & Positiva (sobreestimació) \\
H02\_SN & 2.86 & 3.35 & -14.6 & Sí (+10\%) & Negativa (subestimació) \\
N06\_SN & 4.16 & 1.95 & 113.3 & Sí (+10\%) & Positiva
(sobreestimació) \\
N06A\_SN & 4.58 & 2.09 & 119.1 & Sí (+10\%) & Positiva
(sobreestimació) \\
N06AA\_SN & 1.83 & 0.84 & 117.9 & Sí (+10\%) & Positiva
(sobreestimació) \\
N06AB\_SN & 4.93 & 2.76 & 78.6 & Sí (+10\%) & Positiva
(sobreestimació) \\
N06AX\_SN & 2.98 & 0.96 & 210.4 & Sí (+10\%) & Positiva
(sobreestimació) \\
R03BA\_SN & 1.14 & 1.74 & -34.5 & Sí (+10\%) & Negativa
(subestimació) \\
\end{longtable}

\section{Discussió}\label{discussiuxf3}

\section{Avaluació de l'article original des del punt de vista
estadístic i
metodològic}\label{avaluaciuxf3-de-larticle-original-des-del-punt-de-vista-estaduxedstic-i-metodoluxf2gic}

\clearpage
\appendix

\section{Annex}\label{annex}

\subsection{Creuar data}\label{creuar-data}

\label{anx:carrega}

\begin{Shaded}
\begin{Highlighting}[]
\FunctionTok{library}\NormalTok{(haven)}
\FunctionTok{library}\NormalTok{(dplyr)}

\NormalTok{uab\_demo3 }\OtherTok{\textless{}{-}} \FunctionTok{read\_sas}\NormalTok{(}\StringTok{"uab\_demo3.sas7bdat"}\NormalTok{, }\ConstantTok{NULL}\NormalTok{)}
\NormalTok{uab\_drugs3 }\OtherTok{\textless{}{-}} \FunctionTok{read\_sas}\NormalTok{(}\StringTok{"uab\_drugs3.sas7bdat"}\NormalTok{, }\ConstantTok{NULL}\NormalTok{)}
\NormalTok{uab\_atcdrugs3 }\OtherTok{\textless{}{-}} \FunctionTok{read\_sas}\NormalTok{(}\StringTok{"uab\_atc\_drugs3.sas7bdat"}\NormalTok{, }\ConstantTok{NULL}\NormalTok{)}

\NormalTok{data1 }\OtherTok{\textless{}{-}} \FunctionTok{merge}\NormalTok{(uab\_demo3, uab\_drugs3, }\AttributeTok{by.x=}\StringTok{"PatNo"}\NormalTok{, }\AttributeTok{by.y=}\StringTok{"PatNo"}\NormalTok{)}
\NormalTok{data }\OtherTok{\textless{}{-}} \FunctionTok{merge}\NormalTok{(data1, uab\_atcdrugs3, }\AttributeTok{by.x=}\StringTok{"PatNo"}\NormalTok{, }\AttributeTok{by.y=}\StringTok{"PatNo"}\NormalTok{)}

\NormalTok{data }\OtherTok{\textless{}{-}}\NormalTok{ data }\SpecialCharTok{\%\textgreater{}\%}
  \FunctionTok{mutate}\NormalTok{(}\FunctionTok{across}\NormalTok{(}\FunctionTok{where}\NormalTok{(}\SpecialCharTok{\textasciitilde{}} \FunctionTok{all}\NormalTok{(}\FunctionTok{unique}\NormalTok{(.x) }\SpecialCharTok{\%in\%} \FunctionTok{c}\NormalTok{(}\DecValTok{0}\NormalTok{, }\DecValTok{1}\NormalTok{, }\DecValTok{2}\NormalTok{, }\DecValTok{3}\NormalTok{, }\ConstantTok{NA}\NormalTok{,}\StringTok{"M"}\NormalTok{,}\StringTok{"F"}\NormalTok{))), as.factor))}
\NormalTok{data}\SpecialCharTok{$}\NormalTok{NFractura}\OtherTok{=}\FunctionTok{as.factor}\NormalTok{(data}\SpecialCharTok{$}\NormalTok{NFractura)}
\FunctionTok{levels}\NormalTok{(data}\SpecialCharTok{$}\NormalTok{NFractura)[}\FunctionTok{levels}\NormalTok{(data}\SpecialCharTok{$}\NormalTok{NFractura) }\SpecialCharTok{\%in\%} \FunctionTok{c}\NormalTok{(}\StringTok{"2"}\NormalTok{, }\StringTok{"3"}\NormalTok{, }\StringTok{"4"}\NormalTok{, }\StringTok{"5"}\NormalTok{, }\StringTok{"6"}\NormalTok{, }\StringTok{"7"}\NormalTok{)]}\OtherTok{=}\StringTok{"\textgreater{}=2"}
\NormalTok{data}\SpecialCharTok{$}\NormalTok{ServeiAlta}\OtherTok{=}\FunctionTok{as.factor}\NormalTok{(data}\SpecialCharTok{$}\NormalTok{ServeiAlta)}
\NormalTok{data}\SpecialCharTok{$}\NormalTok{Edat}\OtherTok{=}\FunctionTok{cut}\NormalTok{(data}\SpecialCharTok{$}\NormalTok{Edat,}\AttributeTok{breaks=}\FunctionTok{c}\NormalTok{(}\DecValTok{50}\NormalTok{, }\DecValTok{65}\NormalTok{, }\DecValTok{81}\NormalTok{, }\DecValTok{95}\NormalTok{),}\AttributeTok{labels =} \FunctionTok{c}\NormalTok{(}\StringTok{"51{-}65"}\NormalTok{,}\StringTok{"66{-}81"}\NormalTok{,}\StringTok{"82{-}95"}\NormalTok{),}\AttributeTok{include.lowest =} \ConstantTok{TRUE}\NormalTok{)}
\NormalTok{data}\OtherTok{=}\FunctionTok{subset}\NormalTok{(data, }\AttributeTok{select=}\SpecialCharTok{{-}}\FunctionTok{c}\NormalTok{(Hipogonad, FX\_familia))}
\NormalTok{data}\SpecialCharTok{$}\NormalTok{Smoker}\OtherTok{=}\FunctionTok{factor}\NormalTok{(data}\SpecialCharTok{$}\NormalTok{Smoker,}\AttributeTok{levels =} \FunctionTok{c}\NormalTok{(}\DecValTok{0}\NormalTok{,}\DecValTok{1}\NormalTok{,}\DecValTok{2}\NormalTok{),}\AttributeTok{labels =} \FunctionTok{c}\NormalTok{(}\StringTok{"No Fumador"}\NormalTok{, }\StringTok{"Fumador Actiu"}\NormalTok{, }\StringTok{"Ex{-}Fumador"}\NormalTok{))}
\NormalTok{data}\SpecialCharTok{$}\NormalTok{ALCOHOL}\OtherTok{=}\FunctionTok{factor}\NormalTok{(data}\SpecialCharTok{$}\NormalTok{ALCOHOL,}\AttributeTok{levels =} \FunctionTok{c}\NormalTok{(}\DecValTok{0}\NormalTok{,}\DecValTok{1}\NormalTok{,}\DecValTok{2}\NormalTok{),}\AttributeTok{labels =} \FunctionTok{c}\NormalTok{(}\StringTok{"No consumeix"}\NormalTok{, }\StringTok{"Poc consum"}\NormalTok{,}\StringTok{"Consum moderat"}\NormalTok{))}
\NormalTok{data}\SpecialCharTok{$}\NormalTok{tipus\_diabetis}\OtherTok{=}\FunctionTok{factor}\NormalTok{(data}\SpecialCharTok{$}\NormalTok{tipus\_diabetis,}\AttributeTok{levels =} \FunctionTok{c}\NormalTok{(}\DecValTok{0}\NormalTok{,}\DecValTok{1}\NormalTok{,}\DecValTok{2}\NormalTok{),}\AttributeTok{labels =} \FunctionTok{c}\NormalTok{(}\StringTok{"No diabetis"}\NormalTok{,}\StringTok{"Tipus 1"}\NormalTok{,}\StringTok{"Tipus 2"}\NormalTok{))}
\end{Highlighting}
\end{Shaded}

\subsection{Taula característiques}\label{taula-caracteruxedstiques}

\label{anx:caract}

\begin{Shaded}
\begin{Highlighting}[]
\FunctionTok{library}\NormalTok{(dplyr)}
\FunctionTok{library}\NormalTok{(tableone)}

\NormalTok{variables\_demo }\OtherTok{\textless{}{-}} \FunctionTok{c}\NormalTok{(}\StringTok{"Sexe"}\NormalTok{, }\StringTok{"Edat"}\NormalTok{ ,}\StringTok{"PES"}\NormalTok{, }\StringTok{"TALLA"}\NormalTok{, }\StringTok{"IMC"}\NormalTok{,}
                  \StringTok{"Smoker"}\NormalTok{,}\StringTok{"ALCOHOL"}\NormalTok{, }\StringTok{"Artritis\_reumatoide"}\NormalTok{,}
                  \StringTok{"Fractura"}\NormalTok{,}\StringTok{"NFractura"}\NormalTok{,}\StringTok{"Diabetis"}\NormalTok{,}\StringTok{"tipus\_diabetis"}\NormalTok{,}\StringTok{"Osteporosi"}\NormalTok{, }
                  \StringTok{"Densitometries"}\NormalTok{,}\StringTok{"neoplasia"}\NormalTok{,}\StringTok{"HiperTiroidisme"}\NormalTok{,}\StringTok{"Malnutricio"}\NormalTok{,}
                  \StringTok{"Malabsorcio"}\NormalTok{,}\StringTok{"Malaltia\_Hep\_Cro"}\NormalTok{,}\StringTok{"OthRiskFract"}\NormalTok{, }\StringTok{"CountActualGE4"}\NormalTok{)}

\NormalTok{cat\_vars }\OtherTok{\textless{}{-}} \FunctionTok{c}\NormalTok{(}\StringTok{"Sexe"}\NormalTok{, }\StringTok{"Smoker"}\NormalTok{, }\StringTok{"ALCOHOL"}\NormalTok{, }\StringTok{"Artritis\_reumatoide"}\NormalTok{,}\StringTok{"Edat"}\NormalTok{, }
              \StringTok{"Fractura"}\NormalTok{, }\StringTok{"Diabetis"}\NormalTok{, }\StringTok{"tipus\_diabetis"}\NormalTok{, }\StringTok{"Osteporosi"}\NormalTok{, }\StringTok{"Densitometries"}\NormalTok{,}
              \StringTok{"neoplasia"}\NormalTok{, }\StringTok{"HiperTiroidisme"}\NormalTok{, }\StringTok{"Malnutricio"}\NormalTok{, }\StringTok{"Malabsorcio"}\NormalTok{, }
              \StringTok{"Malaltia\_Hep\_Cro"}\NormalTok{,}\StringTok{"OthRiskFract"}\NormalTok{, }\StringTok{"CountActualGE4"}\NormalTok{)}

\NormalTok{taula\_demo }\OtherTok{\textless{}{-}} \FunctionTok{CreateTableOne}\NormalTok{(}
  \AttributeTok{vars =}\NormalTok{ variables\_demo,}
  \AttributeTok{strata =} \StringTok{"ControlCas.x"}\NormalTok{,}
  \AttributeTok{data =}\NormalTok{ data,}
  \AttributeTok{factorVars =}\NormalTok{ cat\_vars,}
  \AttributeTok{includeNA =} \ConstantTok{FALSE}
\NormalTok{)}

\NormalTok{tbl}\OtherTok{\textless{}{-}} \FunctionTok{print}\NormalTok{(taula\_demo)}
\FunctionTok{colnames}\NormalTok{(tbl)}\OtherTok{=}\FunctionTok{c}\NormalTok{(}\StringTok{"Control"}\NormalTok{,}\StringTok{"Cas"}\NormalTok{,}\StringTok{"P\_valor"}\NormalTok{,}\StringTok{" "}\NormalTok{)}
\end{Highlighting}
\end{Shaded}

\subsection{Anàlisi d'exposició als fàrmacs i els riscs associats de
fractura}\label{anuxe0lisi-dexposiciuxf3-als-fuxe0rmacs-i-els-riscs-associats-de-fractura}

\subsubsection{Anàlisi descriptiva}\label{anuxe0lisi-descriptiva}

\label{anx:descrip}

\begin{Shaded}
\begin{Highlighting}[]
\NormalTok{variables\_farmacs }\OtherTok{\textless{}{-}} \FunctionTok{c}\NormalTok{(}\StringTok{"CortInh"}\NormalTok{,}
\StringTok{"CortSist"}\NormalTok{,}\StringTok{"CortSistIniBf3m"}\NormalTok{,}\StringTok{"CortSistExpLt3m"}\NormalTok{,}\StringTok{"CortInhIniBf3m"}\NormalTok{,}
\StringTok{"CortInhExpLt3m"}\NormalTok{,}\StringTok{"ADO\_no\_glitazona"}\NormalTok{,}\StringTok{"Bisf"}\NormalTok{,}\StringTok{"ADepreISRS"}\NormalTok{,}
\StringTok{"ADepreNoISRS"}\NormalTok{,}\StringTok{"insulina"}\NormalTok{,}\StringTok{"H\_SN"}\NormalTok{,}\StringTok{"N\_SN"}\NormalTok{,}\StringTok{"H01\_SN"}\NormalTok{,}\StringTok{"H02\_SN"}\NormalTok{,}
\StringTok{"N06\_SN"}\NormalTok{,}\StringTok{"N06A\_SN"}\NormalTok{,}\StringTok{"N06AA\_SN"}\NormalTok{,}\StringTok{"N06AB\_SN"}\NormalTok{,}\StringTok{"N06AX\_SN"}\NormalTok{,}\StringTok{"R03BA\_SN"}\NormalTok{)}

\NormalTok{variables\_cat }\OtherTok{\textless{}{-}}\NormalTok{ variables\_farmacs[}\SpecialCharTok{!}\FunctionTok{grepl}\NormalTok{(}\StringTok{"Dias"}\NormalTok{, variables\_farmacs)]}

\NormalTok{taula\_des }\OtherTok{\textless{}{-}} \FunctionTok{CreateTableOne}\NormalTok{(}
  \AttributeTok{vars =}\NormalTok{ variables\_farmacs,}
  \AttributeTok{factorVars =}\NormalTok{ variables\_cat,}
  \AttributeTok{strata =} \StringTok{"ControlCas.x"}\NormalTok{,}
  \AttributeTok{data =}\NormalTok{ data}
\NormalTok{)}

\NormalTok{tbl}\OtherTok{=}\FunctionTok{print}\NormalTok{(taula\_des)}
\FunctionTok{colnames}\NormalTok{(tbl)}\OtherTok{=}\FunctionTok{c}\NormalTok{(}\StringTok{"Control"}\NormalTok{,}\StringTok{"Cas"}\NormalTok{,}\StringTok{"P\_valor"}\NormalTok{,}\StringTok{" "}\NormalTok{)}
\FunctionTok{class}\NormalTok{(tbl)}
\end{Highlighting}
\end{Shaded}

\begin{Shaded}
\begin{Highlighting}[]
\FunctionTok{library}\NormalTok{(knitr)}
\FunctionTok{kable}\NormalTok{(tbl,}\AttributeTok{caption =} \StringTok{"}\SpecialCharTok{\textbackslash{}\textbackslash{}}\StringTok{label\{tab:descrip\}Comparació de Medicaments i teràpies"}\NormalTok{,}\AttributeTok{align =} \StringTok{"lccc"}\NormalTok{)}
\end{Highlighting}
\end{Shaded}

\subsubsection{Anàlisi logística
``crua''}\label{anuxe0lisi-loguxedstica-crua-1}

\begin{Shaded}
\begin{Highlighting}[]
\FunctionTok{library}\NormalTok{(survival)}

\NormalTok{data}\SpecialCharTok{$}\NormalTok{cas }\OtherTok{\textless{}{-}} \FunctionTok{as.numeric}\NormalTok{(}\FunctionTok{as.character}\NormalTok{(data}\SpecialCharTok{$}\NormalTok{ControlCas.x))}
\NormalTok{data}\SpecialCharTok{$}\NormalTok{strata\_id }\OtherTok{\textless{}{-}} \FunctionTok{as.factor}\NormalTok{(data}\SpecialCharTok{$}\NormalTok{MatchCC.x)}

\NormalTok{variables\_a\_excloure }\OtherTok{\textless{}{-}} \FunctionTok{c}\NormalTok{(}\StringTok{"PatNo"}\NormalTok{, }\StringTok{"ServeiAlta"}\NormalTok{, }\StringTok{"Sexe"}\NormalTok{, }\StringTok{"Edat"}\NormalTok{, }\StringTok{"MatchCC"}\NormalTok{, }\StringTok{"MatchCC.x"}\NormalTok{, }\StringTok{"MatchCC.y"}\NormalTok{, }\StringTok{"PES"}\NormalTok{, }\StringTok{"TALLA"}\NormalTok{, }\StringTok{"IMC"}\NormalTok{,}
                  \StringTok{"Smoker"}\NormalTok{,}\StringTok{"ALCOHOL"}\NormalTok{, }\StringTok{"Artritis\_reumatoide"}\NormalTok{,}
                  \StringTok{"Fractura"}\NormalTok{,}\StringTok{"NFractura"}\NormalTok{,}\StringTok{"Diabetis"}\NormalTok{,}\StringTok{"tipus\_diabetis"}\NormalTok{,}\StringTok{"Osteporosi"}\NormalTok{, }
                  \StringTok{"Densitometries"}\NormalTok{,}\StringTok{"neoplasia"}\NormalTok{,}\StringTok{"HiperTiroidisme"}\NormalTok{,}\StringTok{"Malnutricio"}\NormalTok{,}
                  \StringTok{"Malabsorcio"}\NormalTok{,}\StringTok{"Malaltia\_Hep\_Cro"}\NormalTok{,}\StringTok{"OthRiskFract"}\NormalTok{, }\StringTok{"CountActualGE4"}\NormalTok{, }
  \StringTok{"ControlCas.x"}\NormalTok{, }\StringTok{"ControlCas.y"}\NormalTok{, }\StringTok{"ControlCas"}\NormalTok{, }\StringTok{"cas"}\NormalTok{, }\StringTok{"strata\_id"}\NormalTok{)}

\NormalTok{totes\_les\_columnes }\OtherTok{\textless{}{-}} \FunctionTok{names}\NormalTok{(data)}

\NormalTok{variables\_analisi }\OtherTok{\textless{}{-}} \FunctionTok{setdiff}\NormalTok{(totes\_les\_columnes, variables\_a\_excloure)}

\NormalTok{resultats\_2b }\OtherTok{\textless{}{-}} \FunctionTok{lapply}\NormalTok{(variables\_analisi, }\ControlFlowTok{function}\NormalTok{(variable) \{}

\NormalTok{    formula\_model }\OtherTok{\textless{}{-}} \FunctionTok{as.formula}\NormalTok{(}
    \FunctionTok{paste}\NormalTok{(}\StringTok{"cas \textasciitilde{}"}\NormalTok{, variable, }\StringTok{"+ strata(strata\_id)"}\NormalTok{)}
\NormalTok{  )}
  
\NormalTok{  model }\OtherTok{\textless{}{-}} \FunctionTok{clogit}\NormalTok{(formula\_model, }\AttributeTok{data =}\NormalTok{ data)}
\NormalTok{  resum }\OtherTok{\textless{}{-}} \FunctionTok{summary}\NormalTok{(model)}\SpecialCharTok{$}\NormalTok{coefficients}
    
\NormalTok{    coeficients }\OtherTok{\textless{}{-}} \FunctionTok{summary}\NormalTok{(model)}\SpecialCharTok{$}\NormalTok{coefficients}

\NormalTok{    estimacio }\OtherTok{\textless{}{-}}\NormalTok{ estimacio }\OtherTok{\textless{}{-}}\NormalTok{ coeficients[}\DecValTok{1}\NormalTok{]}
\NormalTok{    p\_valor }\OtherTok{\textless{}{-}}\NormalTok{ coeficients[}\DecValTok{5}\NormalTok{]}
    
\NormalTok{    OR }\OtherTok{\textless{}{-}}\NormalTok{ coeficients[}\DecValTok{2}\NormalTok{]}
\NormalTok{    IC\_inf }\OtherTok{\textless{}{-}} \FunctionTok{exp}\NormalTok{(estimacio}\FloatTok{{-}1.96}\SpecialCharTok{*}\NormalTok{coeficients[}\DecValTok{3}\NormalTok{])}
\NormalTok{    IC\_sup }\OtherTok{\textless{}{-}} \FunctionTok{exp}\NormalTok{(estimacio}\FloatTok{+1.96}\SpecialCharTok{*}\NormalTok{coeficients[}\DecValTok{3}\NormalTok{])}
    
    \FunctionTok{return}\NormalTok{(}\FunctionTok{data.frame}\NormalTok{(}
      \AttributeTok{Variable=}\NormalTok{variable,}
      \AttributeTok{OR\_Cru =} \FunctionTok{round}\NormalTok{(OR, }\DecValTok{2}\NormalTok{),}
      \AttributeTok{IC\_Inferior =} \FunctionTok{round}\NormalTok{(IC\_inf, }\DecValTok{2}\NormalTok{),}
      \AttributeTok{IC\_Superior =} \FunctionTok{format}\NormalTok{(}\FunctionTok{round}\NormalTok{(IC\_sup, }\DecValTok{2}\NormalTok{),}\AttributeTok{nsmall=}\DecValTok{2}\NormalTok{),}
      \AttributeTok{P\_valor =} \FunctionTok{round}\NormalTok{(p\_valor, }\DecValTok{3}\NormalTok{)}
\NormalTok{    ))}
\NormalTok{    \})}

\NormalTok{taula\_final\_2b }\OtherTok{\textless{}{-}} \FunctionTok{do.call}\NormalTok{(rbind, resultats\_2b)}
\FunctionTok{rownames}\NormalTok{(taula\_final\_2b) }\OtherTok{\textless{}{-}} \ConstantTok{NULL}

\NormalTok{tab }\OtherTok{\textless{}{-}} \FunctionTok{print}\NormalTok{(taula\_final\_2b)}
\end{Highlighting}
\end{Shaded}

\begin{Shaded}
\begin{Highlighting}[]
\FunctionTok{kable}\NormalTok{(tab, }\AttributeTok{caption=}\StringTok{"}\SpecialCharTok{\textbackslash{}\textbackslash{}}\StringTok{label\{tab:logcru\}Anàlisi per regressió logística crua de les variables"}\NormalTok{,}\AttributeTok{align =} \StringTok{"lccc"}\NormalTok{)}
\end{Highlighting}
\end{Shaded}

\subsubsection{Anàlisi logística
ajustada}\label{anuxe0lisi-loguxedstica-ajustada-1}

\begin{Shaded}
\begin{Highlighting}[]
\NormalTok{covariables }\OtherTok{\textless{}{-}} \FunctionTok{c}\NormalTok{(}\StringTok{"Osteporosi"}\NormalTok{,}\StringTok{"Fractura"}\NormalTok{,}\StringTok{"ALCOHOL"}\NormalTok{,}\StringTok{"Smoker"}\NormalTok{,}\StringTok{"Diabetis"}\NormalTok{,}\StringTok{"CountActualGE4"}\NormalTok{)  }

\NormalTok{resultats\_ajustats }\OtherTok{\textless{}{-}} \FunctionTok{lapply}\NormalTok{(variables\_analisi, }\ControlFlowTok{function}\NormalTok{(variable) \{}

\NormalTok{    formula\_model }\OtherTok{\textless{}{-}} \FunctionTok{as.formula}\NormalTok{(}
        \FunctionTok{paste}\NormalTok{(}\StringTok{"cas \textasciitilde{}"}\NormalTok{, variable, }\StringTok{"+"}\NormalTok{,}\FunctionTok{paste}\NormalTok{(covariables, }\AttributeTok{collapse =} \StringTok{" + "}\NormalTok{),}\StringTok{"+ strata(strata\_id)"}
\NormalTok{        ))  }
    
\NormalTok{    model }\OtherTok{\textless{}{-}} \FunctionTok{clogit}\NormalTok{(formula\_model, }\AttributeTok{data =}\NormalTok{ data)}
\NormalTok{    coeficients }\OtherTok{\textless{}{-}} \FunctionTok{summary}\NormalTok{(model)}\SpecialCharTok{$}\NormalTok{coefficients}
    
\NormalTok{    coeficients }\OtherTok{\textless{}{-}} \FunctionTok{summary}\NormalTok{(model)}\SpecialCharTok{$}\NormalTok{coefficients}
\NormalTok{    estimacio }\OtherTok{\textless{}{-}}\NormalTok{ coeficients[}\DecValTok{1}\NormalTok{,}\DecValTok{1}\NormalTok{]}
\NormalTok{    p\_valor }\OtherTok{\textless{}{-}}\NormalTok{ coeficients[}\DecValTok{1}\NormalTok{,}\DecValTok{5}\NormalTok{]}
\NormalTok{    OR\_aj }\OtherTok{\textless{}{-}}\NormalTok{ coeficients[}\DecValTok{1}\NormalTok{,}\DecValTok{2}\NormalTok{]}
\NormalTok{    IC\_inf }\OtherTok{\textless{}{-}} \FunctionTok{exp}\NormalTok{(estimacio}\FloatTok{{-}1.96}\SpecialCharTok{*}\NormalTok{coeficients[}\DecValTok{1}\NormalTok{,}\DecValTok{3}\NormalTok{])}
\NormalTok{    IC\_sup }\OtherTok{\textless{}{-}} \FunctionTok{exp}\NormalTok{(estimacio}\FloatTok{+1.96}\SpecialCharTok{*}\NormalTok{coeficients[}\DecValTok{1}\NormalTok{,}\DecValTok{3}\NormalTok{])}
      
    \FunctionTok{return}\NormalTok{(}\FunctionTok{data.frame}\NormalTok{(}
      \AttributeTok{Variable =}\NormalTok{ variable,}
      \AttributeTok{OR\_Ajustat =} \FunctionTok{round}\NormalTok{(OR\_aj, }\DecValTok{2}\NormalTok{),}
      \AttributeTok{IC\_Inferior =} \FunctionTok{round}\NormalTok{(IC\_inf, }\DecValTok{2}\NormalTok{),}
      \AttributeTok{IC\_Superior =} \FunctionTok{format}\NormalTok{(}\FunctionTok{round}\NormalTok{(IC\_sup, }\DecValTok{2}\NormalTok{),}\AttributeTok{nsmall=}\DecValTok{2}\NormalTok{),}
      \AttributeTok{P\_valor =} \FunctionTok{round}\NormalTok{(p\_valor, }\DecValTok{3}\NormalTok{)}
\NormalTok{    ))}
\NormalTok{\})}


\NormalTok{taula\_final\_ajustada }\OtherTok{\textless{}{-}} \FunctionTok{do.call}\NormalTok{(rbind, resultats\_ajustats)}
\FunctionTok{rownames}\NormalTok{(taula\_final\_ajustada) }\OtherTok{\textless{}{-}} \ConstantTok{NULL}
\NormalTok{tab2 }\OtherTok{\textless{}{-}} \FunctionTok{print}\NormalTok{(taula\_final\_ajustada)}
\end{Highlighting}
\end{Shaded}

\begin{Shaded}
\begin{Highlighting}[]
\FunctionTok{kable}\NormalTok{(tab2, }\AttributeTok{caption =} \StringTok{"}\SpecialCharTok{\textbackslash{}\textbackslash{}}\StringTok{label\{tab:logajust\}Anàlisi per regressió logística ajustada per covariables"}\NormalTok{,}\AttributeTok{align =} \StringTok{"lcccc"}\NormalTok{)}
\end{Highlighting}
\end{Shaded}

\subsection{Estadístics descriptius, mesures de risc,
inferencies\ldots{}}\label{estaduxedstics-descriptius-mesures-de-risc-inferencies-1}

\begin{Shaded}
\begin{Highlighting}[]
\CommentTok{\#Taula descriptiva}
\NormalTok{taula\_des\_df }\OtherTok{\textless{}{-}} \FunctionTok{data.frame}\NormalTok{(}
  \AttributeTok{Variable =} \FunctionTok{rownames}\NormalTok{(tbl),}
  \AttributeTok{Controls =}\NormalTok{ tbl[, }\StringTok{"Control"}\NormalTok{],}
  \AttributeTok{Casos =}\NormalTok{ tbl[, }\StringTok{"Cas"}\NormalTok{],}
  \AttributeTok{stringsAsFactors =} \ConstantTok{FALSE}
\NormalTok{)}

\NormalTok{taula\_des\_df }\OtherTok{\textless{}{-}}\NormalTok{ taula\_des\_df }\SpecialCharTok{\%\textgreater{}\%}
  \FunctionTok{filter}\NormalTok{(}\FunctionTok{grepl}\NormalTok{(}\StringTok{"= 1 }\SpecialCharTok{\textbackslash{}\textbackslash{}}\StringTok{(\%}\SpecialCharTok{\textbackslash{}\textbackslash{}}\StringTok{)"}\NormalTok{, Variable))}

\NormalTok{taula\_des\_df}\SpecialCharTok{$}\NormalTok{Variable }\OtherTok{\textless{}{-}} \FunctionTok{gsub}\NormalTok{(}\StringTok{" = 1 }\SpecialCharTok{\textbackslash{}\textbackslash{}}\StringTok{(\%}\SpecialCharTok{\textbackslash{}\textbackslash{}}\StringTok{)"}\NormalTok{, }\StringTok{""}\NormalTok{, taula\_des\_df}\SpecialCharTok{$}\NormalTok{Variable)}

\CommentTok{\#OR crua }
\NormalTok{taula\_final\_2b}\SpecialCharTok{$}\NormalTok{Crua }\OtherTok{\textless{}{-}} \FunctionTok{paste0}\NormalTok{(}
\NormalTok{  taula\_final\_2b}\SpecialCharTok{$}\NormalTok{OR\_Cru, }\StringTok{" ("}\NormalTok{,}
\NormalTok{  taula\_final\_2b}\SpecialCharTok{$}\NormalTok{IC\_Inferior, }\StringTok{","}\NormalTok{,}
\NormalTok{  taula\_final\_2b}\SpecialCharTok{$}\NormalTok{IC\_Superior, }\StringTok{"); p="}\NormalTok{,}
\NormalTok{  taula\_final\_2b}\SpecialCharTok{$}\NormalTok{P\_valor}
\NormalTok{)}

\CommentTok{\#OR ajustada}
\NormalTok{taula\_final\_ajustada}\SpecialCharTok{$}\NormalTok{Ajustada }\OtherTok{\textless{}{-}} \FunctionTok{paste0}\NormalTok{(}
\NormalTok{  taula\_final\_ajustada}\SpecialCharTok{$}\NormalTok{OR\_Ajustat, }\StringTok{" ("}\NormalTok{,}
\NormalTok{  taula\_final\_ajustada}\SpecialCharTok{$}\NormalTok{IC\_Inferior, }\StringTok{","}\NormalTok{,}
\NormalTok{  taula\_final\_ajustada}\SpecialCharTok{$}\NormalTok{IC\_Superior, }\StringTok{"); p="}\NormalTok{,}
\NormalTok{  taula\_final\_ajustada}\SpecialCharTok{$}\NormalTok{P\_valor}
\NormalTok{)}

\NormalTok{taula\_final }\OtherTok{\textless{}{-}}\NormalTok{ taula\_des\_df }\SpecialCharTok{\%\textgreater{}\%}
  \FunctionTok{left\_join}\NormalTok{(taula\_final\_2b[, }\FunctionTok{c}\NormalTok{(}\StringTok{"Variable"}\NormalTok{, }\StringTok{"Crua"}\NormalTok{)], }\AttributeTok{by =} \StringTok{"Variable"}\NormalTok{) }\SpecialCharTok{\%\textgreater{}\%}
  \FunctionTok{left\_join}\NormalTok{(taula\_final\_ajustada[, }\FunctionTok{c}\NormalTok{(}\StringTok{"Variable"}\NormalTok{, }\StringTok{"Ajustada"}\NormalTok{)], }\AttributeTok{by =} \StringTok{"Variable"}\NormalTok{)}

\NormalTok{taula\_final }\OtherTok{\textless{}{-}}\NormalTok{ taula\_final }\SpecialCharTok{\%\textgreater{}\%}
  \FunctionTok{select}\NormalTok{(}\AttributeTok{Exposure =}\NormalTok{ Variable,}
         \StringTok{\textasciigrave{}}\AttributeTok{Controls n (\%)}\StringTok{\textasciigrave{}} \OtherTok{=}\NormalTok{ Controls,}
         \StringTok{\textasciigrave{}}\AttributeTok{Casos n (\%)}\StringTok{\textasciigrave{}} \OtherTok{=}\NormalTok{ Casos,}
         \StringTok{\textasciigrave{}}\AttributeTok{OR crua (IC95\%); p valor}\StringTok{\textasciigrave{}} \OtherTok{=}\NormalTok{ Crua,}
         \StringTok{\textasciigrave{}}\AttributeTok{OR ajustada (IC95\%); p valor}\StringTok{\textasciigrave{}} \OtherTok{=}\NormalTok{ Ajustada)}

\NormalTok{tab3 }\OtherTok{\textless{}{-}} \FunctionTok{print}\NormalTok{(taula\_final)}
\end{Highlighting}
\end{Shaded}

\begin{Shaded}
\begin{Highlighting}[]
\FunctionTok{kable}\NormalTok{(taula\_final, }\AttributeTok{caption =} \StringTok{"Exposició a fàrmacs i riscos associats"}\NormalTok{, }\AttributeTok{align =} \StringTok{"lcccc"}\NormalTok{)}
\end{Highlighting}
\end{Shaded}

\clearpage

\subsection*{Referències}\label{referuxe8ncies}
\addcontentsline{toc}{subsection}{Referències}

\phantomsection\label{refs}
\begin{CSLReferences}{1}{0}
\bibitem[\citeproctext]{ref-brotomodulo}
Broto, Pilar Lalueza, Lourdes Girona Brumós, and Ramón Ribera Montañá.
n.d. {``M{Ó}DULO 11. INTERACCIONES FARMACOL{Ó}GICAS DE f{Á}RMACOS CON
ACCI{Ó}n SOBRE EL SISTEMA NERVIOSO CENTRAL: ANSIOL{Í}TICOS e
HIPN{Ó}TICOS, ANTIDEPRESIVOS, ANTIPSIC{Ó}TICOS.''}

\bibitem[\citeproctext]{ref-coupland2011antidepressant}
Coupland, Carol, Paula Dhiman, Richard Morriss, Antony Arthur, Garry
Barton, and Julia Hippisley-Cox. 2011. {``Antidepressant Use and Risk of
Adverse Outcomes in Older People: Population Based Cohort Study.''}
\emph{Bmj} 343.

\bibitem[\citeproctext]{ref-gorgas2021}
Gorgas, Maria Queralt, Ferran Torres, Roser Vives, Irene Lopez-Rico,
Dolors Capella, and Caridad Pontes. 2021. {``Effects of Selective
Serotonin Reuptake Inhibitors and Other Antidepressant Drugs on the Risk
of Hip Fracture: A Case-Control Study in an Elderly Mediterranean
Population.''} \emph{European Journal of Hospital Pharmacy} 28: 28--32.
\url{https://doi.org/10.1136/ejhpharm-2019-001893}.

\end{CSLReferences}

\end{document}
